\section{A second return channel: heterogeneous recovery sensing}
\label{sec:sensed-recovery}

The curvature return in \cref{eq:response-functional} is not essential to
the hidden-return principle.  We now change the recovery law while retaining
the same delayed voltage network.  This produces a second model family in
which the curvature channel can vanish identically and the response is
carried instead by heterogeneous sensing of the transverse voltage.

Let \(\varpi_N,d_N^{\rm rec}\in\R^N\) satisfy the common bounds
\begin{equation}
 0<\varpi_-\le \varpi_{i,N}\le\varpi_+,\qquad
 \pi_N^T\varpi_N=1,\qquad
 \sup_N\|d_N^{\rm rec}\|_\infty\le d_+.              \label{eq:sensing-bounds}
\end{equation}
Retain \cref{eq:shared-rfde-v}, but replace
\cref{eq:shared-rfde-w} by
\begin{equation}
 \dot w=\delta^2\left\{1+
 \pi_N^T\left[\varpi_N\circ(v-\mathbf1)
       +d_N^{\rm rec}\circ(v-\mathbf1)^{\circ2}\right]-a\right\}.
                                                               \label{eq:sensed-recovery-law}
\end{equation}
The original family is the special case
\(\varpi_N=\mathbf1\), \(d_N^{\rm rec}=0\).

Put
\begin{align}
 \chi_N^{\rm rec}
 &=\pi_N^T(\varpi_N\circ g_N+d_N^{\rm rec}),\notag\\
 \nu_{0,N}^{\rm rec}
 &=-\frac{3\kappa_N}{8\alpha^3}
      +\frac{\chi_N^{\rm rec}}{4\alpha^2},           \label{eq:sensing-center}\\
 h_{*,N}(\mathbf R)
 &=\frac{K}{2\alpha}A_N^{-1}P_{\perp,N}
       \left(\sum_k\theta_kR_k\right)\mathbf1,       \label{eq:sensing-transverse-source}\\
 e_N&=\varpi_N-\alpha^{-1}c_N,\qquad \pi_N^Te_N=0,\notag\\
 \Lambda_N^{\rm rec}(\mathbf R)
 &=\pi_N^T\{e_N\circ h_{*,N}(\mathbf R)\}.           \label{eq:sensing-response-functional}
\end{align}
For a fixed admissible preparation applied to this polynomial field, write
\(z_N^{a,{\rm rec}}\) and \(z_N^{r,{\rm rec}}\) for the selected traces and
define their normalized central gap by
\begin{equation}
 \mathscr D_N^{\rm rec}(\delta,\nu,\zeta)
 =\frac{2}{\alpha e}\left\{
   \mathscr H_\alpha(z_N^{a,{\rm rec}}(0))
  -\mathscr H_\alpha(z_N^{r,{\rm rec}}(0))\right\}.  \label{eq:sensing-gap}
\end{equation}

\begin{lemma}[Range and finite-section transfer for sensed recovery]
\label[lemma]{lem:sensed-range-gap-transfer}
Assume
\cref{eq:shared-markov,eq:shared-curvature,eq:shared-layer-bounds,%
eq:sensing-bounds}, fix \(p>4\), and fix a compact interval \(I_\nu\)
containing every \(\nu_{0,N}^{\rm rec}\) with a common positive boundary
distance.  For one fixed admissible preparation constructed from
\cref{eq:sensed-recovery-law}, all conclusions of the uniform graph and
trace range in \cref{lem:shared-graph-trace} hold for sensed-recovery
coefficients, with constants independent of \(N\).  The physical history
lift has a continuous left inverse and is injective; in particular,
\begin{equation}
 \mathscr D_N^{\rm rec}=0
 \quad\Longleftrightarrow\quad
 \text{the two selected complete RFDE states agree}.           \label{eq:sensing-zero-fiber}
\end{equation}
Moreover, after decreasing common \(\delta\)- and \(\zeta\)-windows, the
gap is jointly \(C^2\) in \((\nu,\zeta)\) and satisfies
\begin{align}
 \frac{\mathscr D_N^{\rm rec}}{\delta}
 &=\sqrt{2\pi}(\nu-\nu_{0,N}^{\rm rec})+O(\delta),
                                                        \label{eq:sensing-gap-expansion}\\
 \partial_\nu\mathscr D_N^{\rm rec}
 &=\delta\sqrt{2\pi}+O(\delta^2),                    \label{eq:sensing-gap-nu}\\
 \partial_\zeta\mathscr D_N^{\rm rec}
 &=-\delta^2\sqrt{2\pi}\Lambda_N^{\rm rec}(\mathbf R_N)
       +O(\delta^3+\delta^2|\zeta|),                 \label{eq:sensing-gap-zeta}\\
 \partial_{\zeta\zeta}\mathscr D_N^{\rm rec}
 &=O(\delta^2),                                      \label{eq:sensing-gap-zz}\\
 |\partial_{\nu\nu}\mathscr D_N^{\rm rec}|
   +|\partial_{\nu\zeta}\mathscr D_N^{\rm rec}|
 &\le C\delta^2.                                    \label{eq:sensing-gap-mixed}
\end{align}
Every \(O(\cdot)\)-constant may depend on the fixed preparation and the
common model and sensing bounds, but not on \(N\).
\end{lemma}

\begin{theorem}[Hidden response through heterogeneous recovery sensing]
\label[theorem]{thm:sensed-recovery-response}
Assume the common hypotheses
\cref{eq:shared-markov,eq:shared-curvature,eq:shared-layer-bounds,%
eq:sensing-bounds}.  Fix a compact interval \(I_\nu\) containing every
\(\nu_{0,N}^{\rm rec}\) with a common positive boundary distance, and fix
one admissible preparation, with the definition applied to the polynomial
field \cref{eq:sensed-recovery-law}.  There are
\(\delta_0,\zeta_0,c_0,C>0\), independent of \(N\), such that
\(\mathscr D_N^{\rm rec}/\delta\) has exactly one root
\(\nu_{c,N}^{\rm rec}(\delta,\zeta)\) in
\(|\nu-\nu_{0,N}^{\rm rec}|<c_0\) whenever
\(0<\delta<\delta_0\) and \(|\zeta|<\zeta_0\).  No assertion is made about
roots outside this local window.  At this root the two selected complete
RFDE states agree, and \(\zeta\mapsto\nu_{c,N}^{\rm rec}\) is \(C^2\).
With \(\mu_{c,N}^{\rm rec}=\delta^2\nu_{c,N}^{\rm rec}\),
\begin{equation}
 \left|\mu^{\rm rec}_{c,N}(\delta,\zeta)
       -\mu^{\rm rec}_{c,N}(\delta,0)
       -\delta^3\zeta\Lambda_N^{\rm rec}(\mathbf R_N)\right|
 \le C\{\delta^4|\zeta|+\delta^3\zeta^2\},          \label{eq:sensing-root-response}
\end{equation}
with one constant and one parameter box for all \(N\).  For every uniformly
bounded direction family and every fixed admissible preparation, the
baseline-subtracted roots define the response germ
\begin{equation}
 [\mu^{\mathbf p,\rm rec}_{c,N}(\delta,\zeta)
       -\mu^{\mathbf p,\rm rec}_{c,N}(\delta,0)]_{\rm resp}
 =[\delta^3\zeta\Lambda_N^{\rm rec}(\mathbf R_N)]_{\rm resp}. \label{eq:sensing-canonical-germ}
\end{equation}
Any two fixed admissible preparations give the same class after restricting
their representatives to the intersection of their parameter boxes.

For a fixed \(N\), if at least one delay is positive, the restriction of
\(\Lambda_N^{\rm rec}\) to \(\mathfrak T_N^{\rm sr}\) is nonzero if and
only if
\begin{equation}
 \varpi_N\ne \alpha^{-1}c_N.                         \label{eq:sensing-classification}
\end{equation}
If two delays are distinct, the same equivalence holds on
\(\mathfrak T_N^{\rm red}\).  Whenever nonzero, the zero-response
directions form a closed codimension-one hyperplane.
\end{theorem}

\begin{corollary}[Finite-scale recovery of recovery-return pairings]
\label[corollary]{cor:recovery-tomography}
Assume the common network, curvature, sensing, base-layer, and interval
hypotheses of \cref{thm:sensed-recovery-response}, let \(m\ge1\), suppose
\(P_N\) is irreducible, and fix one admissible preparation rule common to the
unit generator-probe family
\(\{\mathsf G_N(d_N):\|d_N\|_N\le1\}\).  Let
\(e_N=\varpi_N-\alpha^{-1}c_N\in E_N\).  For the generator-supported probes
\(\mathsf G_N(d)\) in \cref{eq:tomography-generator-probe},
\begin{equation}
 \Lambda_N^{\rm rec}(\mathsf G_N(d))
 =-\frac{K\Delta_\theta}{2\alpha D}
       \langle e_N,d\rangle_{\pi_N}.                 \label{eq:recovery-tomography-pairing}
\end{equation}
Consequently, responses on any \(N-1\) probes spanning \(E_N\) determine
the complete recovery-return profile \(e_N\).  There are
\(\delta_{\rm tom}^{\rm rec},\zeta_{\rm tom}^{\rm rec},
C_{\rm tom}^{\rm rec}>0\), depending only on the common data in
\cref{thm:sensed-recovery-response} and that fixed preparation rule, but
independent of \(N\) and of the choice of
\(d_N\in E_N\) with \(\|d_N\|_N\le1\), such that the selected roots for
these probes satisfy, for every \(N\ge2\),
\(0<\delta<\delta_{\rm tom}^{\rm rec}\), and
\(0<|\zeta|<\zeta_{\rm tom}^{\rm rec}\),
\begin{equation}
 \left|
  -\frac{2\alpha D}{K\Delta_\theta}
  \frac{\mu^{\mathbf p,d,\rm rec}_{c,N}(\delta,\zeta)
        -\mu^{\mathbf p,d,\rm rec}_{c,N}(\delta,0)}
       {\delta^3\zeta}
  -\langle e_N,d_N\rangle_{\pi_N}
 \right|
 \le C_{\rm tom}^{\rm rec}(\delta+|\zeta|).          \label{eq:recovery-finite-tomography}
\end{equation}

The estimate is uniform for the weighted response pairings.  It does not,
without a uniform lower frame bound for the chosen probes, assert an
\(N\)-uniform condition number for recovering the coordinate vector \(e_N\)
from finitely many noisy pairings.

For each fixed \(N\), varying the sensing weights in the sensed-recovery
model class realizes every nonzero ray in \(E_N^*\).  The probes identify
the fixed ray of a given model; they do not create it.  Indeed, let
\(\mathfrak r\in E_N^*\setminus\{0\}\), and let \(e\in E_N\) be its unique
weighted representative,
\(\mathfrak r(h)=\langle e,h\rangle_{\pi_N}\).  With homogeneous curvature
\(c_N=\alpha\mathbf1\), the choice
\(\varpi_N=\mathbf1+t e\) is admissible for all sufficiently small
\(t>0\) and has recovery-return covector \(t\mathfrak r\).  The same
realization is uniform for families whose representatives have a common
\(\ell^\infty\) bound.
\end{corollary}

\begin{proof}[Proof of \cref{lem:sensed-range-gap-transfer}]
Write \(Z_N=g_NX^2+h\) in the fold chart.  Direct substitution into
\cref{eq:sensed-recovery-law} gives the exact equation
\begin{align}
 Y'={}&-X+\delta\{\nu-\pi_N^T(\varpi_N\circ Z_N)
                    -(\pi_N^Td_N^{\rm rec})X^2\}\notag\\
 &-2\delta^2X\pi_N^T(d_N^{\rm rec}\circ Z_N)
   -\delta^3\pi_N^T(d_N^{\rm rec}\circ Z_N^{\circ2}).
                                                               \label{eq:sensing-exact-Y}
\end{align}
Equivalently, the only change to the prepared normal form
\cref{eq:invdet-raw-normal-form} is
\begin{align}
 F_N^{Y,\natural}={}&\nu-\pi_N^T(\varpi_N\circ W_N)
   -(\pi_N^Td_N^{\rm rec})X_0^2\notag\\
 &-2\rho X_0\pi_N^T(d_N^{\rm rec}\circ W_N)
   -\rho^2\pi_N^T(d_N^{\rm rec}\circ W_N^{\circ2}). \label{eq:sensing-raw-FY}
\end{align}
This is a fixed-degree polynomial with the same dimension-uniform product
bounds as \cref{eq:invdet-raw-F,eq:invdet-raw-G}.  Multiplication by
\(\varpi_N\) and \(d_N^{\rm rec}\) is uniformly bounded in the stationary
oscillation norm by \cref{eq:shared-products,eq:sensing-bounds}, and
\(F_N^{Y,\natural}\) has no direct \(\zeta\)-derivative.  The transfer of
the range construction is therefore not an appeal to the old
model-specific \(Y\)-integrand; its inputs and outputs are as follows.

\begin{center}
\small
\begin{tabular}{@{}L{.18\linewidth}L{.38\linewidth}L{.34\linewidth}@{}}
\toprule
Step & Sensed-recovery input & Output \\
\midrule
Graph contraction &
The singular field, \(A_N\), fixed delay atoms, and semigroup bound are
unchanged; the new prepared component has the same tame derivative
majorant. &
The special-flow transform contracts on the same graph ball.\\
Mixed graph jets &
Each highest block retains the external dummy amplitude \(\rho\);
\(F_N^{Y,\natural}\) is affine in \(\nu\) and directly
\(\zeta\)-independent. &
The diagonal and prolonged inverses, coefficient envelopes, and
\(O(\delta^2)\) structural jets have the same uniform bounds.\\
Trace range &
The singular field, Green operator, phase row, and endpoint row are
unchanged. &
The selected traces and their \(C^2\) parameter jets obey
\cref{eq:shared-trace-jets} and its mixed-derivative extension.\\
Physical lift &
The voltage history determines \(X,h,Z_N\), while the new \(Y\)-equation
does not contain \(Y\). &
A direct integral reconstructs the discarded resource history.\\
Finite gap &
The two new physical integrands are Gaussian times polynomials of degrees
at most four and two. &
The endpoint, tail, moving-trace, parameter, graph-jet, and preparation
errors retain their uniform orders.\\
\bottomrule
\end{tabular}
\end{center}

For the first three rows, insert \cref{eq:sensing-raw-FY} into the
fixed-point and prolonged equations in
\cref{lem:invdet-C0-contraction,prop:invdet-mixed-jets}.  Every
current-block term still carries \(\rho\), while the new multiplication
operators have common norms.  Hence the contraction factor and diagonal
inverse bounds are unchanged after enlarging fixed constants.  The local
block induction in \cref{prop:invdet-coefficient-envelope} gives the same
retained-tube coefficient envelopes.  Since the singular field remains
\(q_0\), the trace contraction and its differentiated equations then give
the range and trace conclusions stated in the lemma.

We also verify the complete-history step rather than inherit it from the old
resource equation.  From the voltage history in the physical RFDE phase
space recover, for \(-\theta_m\le\vartheta\le0\),
\begin{align}
 X(\vartheta)
 &=\delta^{-1}\{\pi_N^Tv(\vartheta/\delta)-1\},\notag\\
 h(\vartheta)
 &=\delta^{-2}P_{\perp,N}v(\vartheta/\delta)
      -g_NX(\vartheta)^2,
 \qquad Z_N(\vartheta)=g_NX(\vartheta)^2+h(\vartheta). \label{eq:sensing-recover-Xh}
\end{align}
The current resource gives \(Y(0)=\delta^{-2}(2/3-w(0))\).  If
\(F_N^{Y,{\rm rec}}(X,Z)\) denotes the right-hand side of
\cref{eq:sensing-exact-Y}, then the discarded resource history is recovered
exactly by
\begin{equation}
 Y(\vartheta)=Y(0)-\int_\vartheta^0
   F_N^{Y,{\rm rec}}(X(r),Z_N(r))\,\dd r.             \label{eq:sensing-recover-Y-history}
\end{equation}
All operations in \cref{eq:sensing-recover-Xh,eq:sensing-recover-Y-history}
are continuous in the physical history norm on the fixed target.  They give
a continuous left inverse of the sensed-recovery history lift.  Hence that
lift is injective, and a zero of \cref{eq:sensing-gap} is equality of the
two complete RFDE states, not merely equality after projection.

The first coefficient and the structural derivative of the second
coefficient on the physical singular orbit are now
\begin{align}
 q^{\rm rec}_{1,N,X}
 &=-\kappa_NX^3+K\pi_N^T\mathcal L_{N,0}[\mathbf1X],\notag\\
 q^{\rm rec}_{1,N,Y}
 &=\nu-\chi_N^{\rm rec}X^2,                          \label{eq:sensing-q1}\\
 D_\zeta q^{\rm rec}_{2,N}(\gamma_0(s),\nu,0)
 &=\begin{pmatrix}
  \dfrac{s}{\alpha}\pi_N^T(c_N\circ h_{*,N})\\[1mm]
  -\pi_N^T(\varpi_N\circ h_{*,N})
 \end{pmatrix}.                                      \label{eq:sensing-q2-return}
\end{align}
Indeed, the first component is the curvature return already computed in
\cref{eq:q2-return}; the second is the derivative of
\(-\delta\pi_N^T(\varpi_N\circ h)\) through the first graph coefficient.
The terms containing \(d_N^{\rm rec}\) change the baseline coefficients but
have no \(\zeta\)-derivative at this order.

With \(\ell_N=-(K/(2\alpha))\pi_N^TM_{1,N}(0)\mathbf1\), the two
whole-line physical integrands become
\begin{align}
 a^{\rm rec}_{1,N}(s,\nu)
 &=e^{-s^2/2}\left\{\nu+
    \frac{\kappa_Ns^4}{8\alpha^3}
    -\frac{\chi_N^{\rm rec}s^2}{4\alpha^2}+\ell_Ns\right\},
                                                               \label{eq:sensing-a1}\\
 a^{\rm rec}_{2,\zeta,N}(s)
 &=e^{-s^2/2}\left\{
   \frac{\pi_N^T(c_N\circ h_{*,N})}{\alpha}s^2
       -\pi_N^T(\varpi_N\circ h_{*,N})\right\}.     \label{eq:sensing-a2}
\end{align}
Gaussian integration gives
\begin{align}
 \int_\R a^{\rm rec}_{1,N}(s,\nu)\,\dd s
 &=\sqrt{2\pi}(\nu-\nu_{0,N}^{\rm rec}),\notag\\
 \int_\R a^{\rm rec}_{2,\zeta,N}(s)\,\dd s
 &=-\sqrt{2\pi}\Lambda_N^{\rm rec}(\mathbf R).     \label{eq:sensing-melnikov}
\end{align}

We now form the finite-section bridge with these two physical integrands.
Define \(A^{\rm rec}_{\ell,N}\), \(A^{\rm rec}_{3,N}\), and the six error
classes by the formulas in
\cref{eq:gap-exact-decomposition,eq:gap-endpoint-error,eq:gap-singular-error,eq:gap-moving-error,eq:gap-parameter-error,eq:gap-jet-error,eq:gap-preparation-error},
replacing \(q_{\ell,N}\), \(\mathcal R_{3,N}\), and the two physical
integrands by their sensed-recovery counterparts.  This is again an exact
finite-section identity.  The endpoint error is zero.  The moving-trace,
parameter-Taylor, graph-jet, and preparation estimates use the trace and
coefficient envelopes just established and are unchanged term by term.
For the singular tails, \cref{eq:sensing-a1} is a Gaussian times a
polynomial of degree at most four and \cref{eq:sensing-a2} one of degree at
most two, with common coefficient bounds.  Thus
\cref{eq:gap-gaussian-tail-bound} applies directly.  The six displayed
error definitions in \cref{prop:finite-gap-bridge} are therefore valid with
the sensed coefficients, and the proof of each bound uses exactly the input
recorded in the last row of the table.  Inserting
\cref{eq:sensing-melnikov} into that exact decomposition gives the five
estimates in \cref{lem:sensed-range-gap-transfer}, with constants and
parameter boxes common in \(N\).
\end{proof}

\begin{proof}[Proof of \cref{thm:sensed-recovery-response}]
The recovery equation has no explicit \(\zeta\)-dependence, so the
stationary-projected right-hand side is exactly blind to \(\zeta\) at every
full history, as in \cref{prop:projection-blindness}.  By
\cref{lem:sensed-range-gap-transfer}, the matcher has an injective
complete-history range, the zero-fiber property, and the five uniform
range-to-gap estimates, with
\[
 a_N=\sqrt{2\pi},\qquad
 b_N=-\sqrt{2\pi}\Lambda_N^{\rm rec}(\mathbf R_N).
\]
Applying \cref{thm:abstract-root-response} proves the unique local root,
its regularity, and \cref{eq:sensing-root-response}.  Applying the same
estimate to any two fixed admissible preparations proves
\cref{eq:sensing-canonical-germ} on their common parameter box.

It remains to classify the covector.  Since \(\pi_N^Te_N=0\), it vanishes
on all transverse responses exactly when \(e_N=0\).  Conversely, if
\(e_N\ne0\), choose a positive delay \(\theta_j\) and set
\begin{equation}
 R_j=\theta_j^{-1}(A_Ne_N)\pi_N^T,
 \qquad R_k=0\quad(k\ne j).                          \label{eq:sensing-single-layer-witness}
\end{equation}
This direction is stationary-row neutral because \(\pi_N^TA_N=0\), and it
satisfies \(R_j\mathbf1=\theta_j^{-1}A_Ne_N\); hence
\(h_{*,N}=(K/(2\alpha))e_N\), and
\(\Lambda_N^{\rm rec}=(K/(2\alpha))\pi_N^T(e_N^{\circ2})\ne0\).
For distinct \(\theta_j,\theta_\ell\), instead set
\begin{equation}
 R_j=(\theta_j-\theta_\ell)^{-1}(A_Ne_N)\pi_N^T,
 \qquad R_\ell=-R_j,                                 \label{eq:sensing-two-layer-witness}
\end{equation}
with all other layers zero.  This gives the same \(h_{*,N}\) and belongs to
\(\mathfrak T_N^{\rm red}\).  The hyperplane assertion follows from
linearity.
\end{proof}

\begin{proof}[Proof of \cref{cor:recovery-tomography}]
The compressed return covector of this model is
\(\mathfrak r_N^{\rm rec}(h)=\langle e_N,h\rangle_{\pi_N}\).
Combining this identity with
\cref{eq:tomography-generator-identity} proves
\cref{eq:recovery-tomography-pairing}.  The weighted pairing is
nondegenerate on \(E_N\), so a spanning probe family uniquely determines
\(e_N\).  The unit probe ball has a common \(\mathfrak T_N\)-bound by
\cref{eq:tomography-generator-condition}.  Apply the sensed-recovery
graph--root proof over the enlarged index set
\(\{(N,d):N\ge2,\ d\in E_N,\ \|d\|_N\le1\}\).  Its model data, fixed
preparation rule, and direction bounds are uniform there, so one parameter
box and one remainder constant work simultaneously.  Substitute
\cref{eq:recovery-tomography-pairing} into
\cref{eq:sensing-root-response} and divide by \(\delta^3\zeta\) to obtain
\cref{eq:recovery-finite-tomography} with one constant for that ball.

Every functional on \(E_N\) has a unique representative in \(E_N\) under
the positive weighted pairing: constants form the annihilator before the
mean-zero normalization is imposed.  If \(c_N=\alpha\mathbf1\) and
\(\varpi_N=\mathbf1+t e\), then
\(\pi_N^T\varpi_N=1\) and the return covector is \(t\mathfrak r\).
Choosing \(t\|e\|_\infty<1\) gives strict positivity.  A common
\(\ell^\infty\) bound permits the same \(t\) for an entire network family.
\end{proof}

The two return mechanisms can be separated exactly.  In the old limit
\(\varpi_N=\mathbf1\), transversality of \(h_{*,N}\) removes the sensing
term and \cref{eq:sensing-response-functional} reduces to
\cref{eq:response-functional}.  In the opposite limit
\(c_N=\alpha\mathbf1\), the curvature term vanishes for every direction,
but heterogeneous \(\varpi_N\) can leave a nonzero response.

\begin{corollary}[Sparse pure-sensing response]
\label[corollary]{cor:pure-sensing-response}
Use the reset--cycle network \cref{eq:reset-cycle-network} and the centered
vector \(s_N\) from \cref{eq:reset-cycle-curvature}.  Set
\[
 \alpha=1,\qquad c_N=\mathbf1,\qquad
 \varpi_N=\mathbf1+\sigma s_N,\qquad d_N^{\rm rec}=0,\qquad
 0<\sigma<(2+\eta)^{-1},\qquad \lambda\ne0.
\]
For \(0=\theta_0<\theta_1=\theta\), take
\begin{equation}
 B_{0,N}=B_{1,N}=\frac{\Id+P_N}{2},\qquad
 R_{0,N}=\lambda(P_N-\Id)\diag(\varpi_N),\qquad
 R_{1,N}=-R_{0,N}.                                   \label{eq:pure-sensing-layers}
\end{equation}
Then the curvature-return covector \(\Lambda_N\) in
\cref{eq:response-functional} is identically zero, whereas
\begin{align}
 \Lambda_N^{\rm rec}(\mathbf R_N)
 &=-\frac{K\lambda\theta}{2D}
       \{\pi_N^T(\varpi_N^{\circ2})-1\},             \label{eq:pure-sensing-coefficient}\\
 \inf_{N\ge3}|\Lambda_N^{\rm rec}(\mathbf R_N)|
 &\ge\frac{2|K\lambda|\theta\sigma^2q^2(1-q)}{D}>0. \label{eq:pure-sensing-lower}
\end{align}
The directions have rank \(N-1\), retain the sparse support of the chosen
base layers,
obey a common operator bound
\begin{equation}
 \sup_{N\ge3}\bigl(\|R_{0,N}\|_{\mathcal L(\R^N,\|\cdot\|_N)}
              +\|R_{1,N}\|_{\mathcal L(\R^N,\|\cdot\|_N)}\bigr)<\infty,
                                                               \label{eq:pure-sensing-uniform-layer-bound}
\end{equation}
and both perturbed layers are entrywise nonnegative whenever
\(2|\zeta\lambda|\|\varpi_N\|_\infty<1\).
The common hypotheses of \cref{thm:sensed-recovery-response} hold.  Since
\(\nu_{0,N}^{\rm rec}=-\beta/8\), a common interval and any fixed admissible
preparation give the nonlinear root response with constants independent of
\(N\).
\end{corollary}

\begin{proof}
For generator-supported directions with
\(Q_N(p)=(P_N-\Id)\diag(p)\), direct use of
\(A_N=D(P_N-\Id)|_{E_N}\) in
\cref{eq:sensing-transverse-source} gives
\begin{equation}
 \Lambda_N^{\rm rec}
 =-\frac{K\lambda(\theta_1-\theta_0)}{2\alpha D}
       \operatorname{Cov}_{\pi_N}(e_N,p).            \label{eq:sensing-generator-covariance}
\end{equation}
Here \(e_N=\sigma s_N\) and \(p=\varpi_N\), so the covariance is
\(\pi_N^T(\varpi_N^{\circ2})-1=\sigma^2\pi_N^T(s_N^{\circ2})\).
The reset--cycle estimate used in \cref{eq:reset-cycle-lower} gives
\(\pi_N^T(s_N^{\circ2})\ge4q^2(1-q)\), proving
\cref{eq:pure-sensing-lower}.  Positivity, support preservation, and rank
follow exactly as in \cref{prop:generator-supported-directions}, since
\(\varpi_N\) is strictly positive.  Finally,
\(\|P_N-\Id\|_{\mathcal L(\R^N,\|\cdot\|_N)}\le2\), while
\cref{eq:shared-products} and the common bound on \(\varpi_N\) give a
uniform norm for multiplication by \(\varpi_N\).  This proves
\cref{eq:pure-sensing-uniform-layer-bound}.
\end{proof}
